\SourceLeaf{432}

also
\[
u=xyz
\]
Es ist folglich
\[
\begin{aligned}F(a+b)&=Fa+Fb\\[7pt]&\quad{}-\sin\operatorname{lemn}a.\sin\operatorname{lemn}b.\sin\operatorname{lemn}(a+b)\end{aligned}
\]
Setzt man
\[
Fa-\frac{a}{90^0}F90^0=Ga
\]
so ist
\[
\begin{aligned}G(a+b)&=Ga+Gb\\[7pt]&\quad{}-\sin\operatorname{lemn}a.\sin\operatorname{lemn}b.\sin\operatorname{lemn}(a+b)\end{aligned}
\]
\[
Ga=\frac{\mathrm dQa}{Qa.\mathrm da}
\]

\CurvaArticle{15}{[15.]}
\[
\begin{aligned}\frac{\varpi}{\pi}&=\Bigl\{1+\Bigl(\frac12\Bigr)^2.\frac12+
\Bigl(\frac{1.3}{2.4}\Bigr)^2.\frac14+
\\[7pt]&\quad{}\Bigl(\frac{1.3.5}{2.4.6}\Bigr)^2\frac18+..\Bigr\}\surd {\frac12}\end{aligned}
\]
\[
\begin{aligned}\frac{\varpi}{\pi}&=\Bigl(1+\frac{1.3}{4.4}.\frac19+
\frac{1.3.5.7}{4.4.8.8}.\frac1{81}+
\\[7pt]&\quad{}\frac{1.3.5.7.9.11}{4.4.8.8.12.12}.\frac1{729}+..\Bigr)\surd {\frac23}\end{aligned}
\]
\[
\begin{aligned}\frac{\varpi}{\pi}-\frac2{\varpi}&=\Bigl(\frac14.\frac13+\frac{1.3.5}{4.4.8}.\frac1{27}+
\\[7pt]&\quad{}\frac{1.3.5.7.9}{4.4.8.8.12}.\frac1{243}+
\frac{1.3.5.7.9.11.13}{4.4.8.8.12.12.16}.\frac1{2187}+..\Bigr)\surd {\frac23}\end{aligned}
\]

\CurvaArticle{16}{[16.]}
Ponendo
\[
\frac1{\mathrm M(1,\cos\varphi)}=N+a\cos2\varphi+b\cos4\varphi+\ .\ .\ .
\]
erit
\begingroup\small\[
\begin{array}{rl@{\hspace{12pt}}rl}
N&=1.393203&-\frac1\pi\log(1-\cos2\varphi)&=+0.220635\\
a&=0.581803&&+0.636620\cos2\varphi\\
b&=0.309601&&+0.318310\cos4\varphi\\
c&=0.209449&&+0.212207\cos6\varphi\\
d&=0.157960&&+0.159155\cos8\varphi\\
e&=0.126704.&&+0.127324\cos10\varphi
\end{array}
\]\endgroup

\SourceLeaf{433}

\[
N=\left(\frac{3.3.7.7.11.11.15}{2.4.6.8.10.12.14}\ .\ .\ .\right)^2
=2\frac{\varpi\varpi}{\pi\pi}
\]
\[
a=\frac12\left(\frac{5.5}{4.6}.\frac{9.9}{8.10}.\frac{13.13}{12.14}\ .\ .\ .\right)^2
=\frac4{\varpi\varpi}
\]
\[
b=\frac29N
\]
\[
\frac{\varpi}{\surd 8}=\frac{4.6.8.10.12.14}{5.5.9.9.13.13}\ .\ .\ .
\]
\[
\frac{\varpi}{2}=\frac32.\frac45.\frac76.\frac89.\frac{11}{10}.\frac{12}{13}\ .\ .\ .
\]

\CurvaArticle{17}{[17.]}
\[
\begin{aligned}
&1+\Bigl(\frac12\Bigr)^3x^4+\Bigl(\frac12.\frac34\Bigr)^3x^8+
\\[7pt]&\quad{}\Bigl(\frac12.\frac34.\frac56\Bigr)^3x^{12}+\ .\ .\ .\\[5pt]
&=\Bigl\{1+\Bigl(\frac14\Bigr)^2x^4+\Bigl(\frac14.\frac58\Bigr)^2x^8+
\\[7pt]&\quad{}\Bigl(\frac14.\frac58.\frac9{12}\Bigr)^2x^{12}+\ .\ .\Bigr\}^2
\end{aligned}
\]
\textit{Demonstratio.} Ponatur
\[
1+\left(\frac14\right)^2x^4+\left(\frac14.\frac58\right)^2x^8+\ .\ .\ .=t
\]
eritque, posito
\[
x(x^4-1)\frac{\mathrm{dd}t}{\mathrm dx^2}-(3x^4-1)\frac{\mathrm dt}{\mathrm dx}+x^3t=R,\quad R=0
\]
Hinc etiam
\[
\begin{aligned}&0=x\frac{\mathrm dR}{\mathrm dx}+R;\\[7pt]&\text{nec non}\quad
0=2xt\frac{\mathrm dR}{\mathrm dx}+2\left(t+3x\frac{\mathrm dt}{\mathrm dx}\right)R\end{aligned}
\]
unde fit, evolutione facta
\[
\begin{aligned}
0={}&xx(x^4-1)\left(2t\frac{\mathrm d^3t}{\mathrm dx^3}+6\frac{\mathrm dt}{\mathrm dx}\frac{\mathrm{dd}t}{\mathrm dx^2}\right)
\\[7pt]&\quad{}+3x(3x^4-1)\left(2t\frac{\mathrm{dd}t}{\mathrm dx^2}+2\frac{\mathrm dt}{\mathrm dx}\frac{\mathrm dt}{\mathrm dx}\right)\\[5pt]
&+(19x^4-1)2t\frac{\mathrm dt}{\mathrm dx}+8x^3tt
\end{aligned}
\]
sive ponendo $tt=u$
\[
\begin{aligned}0&=xx(x^4-1)\frac{\mathrm d^3u}{\mathrm dx^3}+3x(3x^4-1)\frac{\mathrm{dd}u}{\mathrm dx^2}
\\[7pt]&\quad{}+(19x^4-1)\frac{\mathrm du}{\mathrm dx}+8x^3u\end{aligned}
\]
cui aequationi invenitur respondere
\[
u=1+\left(\frac12\right)^3x^4+\left(\frac{1.3}{2.4}\right)^3x^8+\ .\ .
\]

\SourceLeaf{434}

Iam quum
\[
t=\frac{\surd 8}{\pi}\int\frac{\mathrm dz}{{}^{4}\!\surd{(1-z^4)(1-x^4z^4)}}
\]
$z=0$ usque ad $z=1$, fit, pro $x=1$,
\[
t=\frac{\surd 8}{\pi}.\frac{\varpi}{2}=\frac{\varpi}{\pi}\surd 2
\]
adeoque
\[
1+\left(\frac12\right)^3+\left(\frac{1.3}{2.4}\right)^3+
\left(\frac{1.3.5}{2.4.6}\right)^3+\ .\ .\ .=2\frac{\varpi\varpi}{\pi\pi}
\]
\begin{center}\textit{Idem alio modo}\end{center}
Valor seriei
\[
1+\left(\frac12\right)^3+\left(\frac{1.3}{2.4}\right)^3+
\left(\frac{1.3.5}{2.4.6}\right)^3+\ .\ .
\]
fit
\[
=\frac4{\pi\pi}\iint\frac{\mathrm d\varphi\,\mathrm d\psi}{\surd {(1-\cos\varphi^2\cos\psi^2)}}
\]
a $\varphi=0$ usque ad $\varphi=90^0$ et a $\psi=0$ usque ad $\psi=90^0$.
Faciendo itaque $\cos\varphi\cos\psi=\cos\upsilon$, idem valor fit
\[
=\frac4{\pi\pi}\iint\frac{\mathrm d\varphi\,\mathrm d\upsilon}{\surd {(\cos\varphi^2-\cos\upsilon^2)}}
\]
et quidem ab $\upsilon=0$ usque ad $\upsilon=90^0$ et a $\varphi=0$ usque ad $\varphi=\upsilon$.
Denique statuendo $\varphi+\upsilon=f$, $\upsilon-\varphi=g$, erit expressio nostra
\[
=\frac2{\pi\pi}\iint\frac{\mathrm df.\mathrm dg}{\surd {\sin f\sin g}}
\]
ab $g=0$ usque ad $g=90^0$ et a $f=g$ usque ad $f=180^0-g$.
Sed haud difficulter probatur integrale eundem valorem nancisci, si sumatur
ab $f=0$ usque ad $f=90^0$ et a $g=0$ usque ad $g=90^0$, unde ipsius valor deducitur
\[
=2\frac{\varpi\varpi}{\pi\pi}
\]
uti supra.

